\section{Приложение к реализуемой сигмоидной регрессии}
\label{reg:sigmoid}
Квазар-выпуклость для реализуемых обобщённых
линейных моделей известна \cite[лемма 3]{regglm}; связь через отношение
производных также используется в \cite[лемма 5]{regfoster}.
Сигмоидные наименьшие квадраты на $\ell_1$-шаре уже применяются
для проверки методов условного градиента \cite[раздел 5.1.4]{khademi}.
Ниже даны явные константы для полученных сертификатов.

Пусть строки $A\in\mathbb R^{n\times d}$ равны $a_i^\top$,
$Q=\{x:\|x\|_1\le R\}$, $R>0$, и $x^*\in Q$.
Для $\sigma(z)=(1+e^{-z})^{-1}$ и ответов
$y_i=\sigma(a_i^\top x^*)$ положим
\begin{equation}\label{reg:sigmoidloss}
 f(x)=\frac1{2n}\sum_{i=1}^n[\sigma(a_i^\top x)-y_i]^2.
\end{equation}
Тогда $f^*=0$ и $\nabla f(x^*)=0$. Обозначим
\[
 \begin{gathered}
 Z=R\max_i\|a_i\|_\infty,\qquad
 m_\sigma=(2+2\cosh Z)^{-1},\quad
 M_\sigma=1/4,\quad C_\sigma=1/(6\sqrt3),\\
 S=A^\top A/n,\quad \lambda=\lambda_{\min}(S),\quad
 \Lambda=\lambda_{\max}(S),\quad K_A=\max_i\|a_i\|^2.
 \end{gathered}
\]
Все допустимые скалярные аргументы лежат в $[-Z,Z]$, где
$m_\sigma\le\sigma'\le M_\sigma$; на всей прямой
$|\sigma''|\le C_\sigma$.

\begin{proposition}[Константы регулярности и кривизны]
\label{reg:sigmoidbounds}
При $\lambda>0$ функция удовлетворяет сдвинутому RSI с
$m=m_\sigma^2\lambda$. Для каждого $t>0$ граница нормы гессиана
выполнена с
\begin{equation}\label{reg:sighess}
 H_0=\Lambda\left(M_\sigma^2+\frac{C_\sigma^2}{2t}\right),
 \qquad H_1=K_A t.
\end{equation}
Кроме того,
\begin{equation}\label{reg:quasar}
 \langle\nabla f(x),x-x^*\rangle\ge\gamma f(x),\qquad
 \gamma=\min\{1,2m_\sigma/M_\sigma\}>0,
\end{equation}
и на всей допустимой области
\begin{equation}\label{reg:sigglobal}
 \|\nabla^2f(x)\|\le L_Q
 :=\Lambda\bigl(M_\sigma^2+C_\sigma\tanh(Z/2)\bigr).
\end{equation}
Для последних двух выводов полный столбцовый ранг не требуется.
\end{proposition}
\begin{proof}
Положим $d=x-x^*$, $z_i=a_i^\top x$ и
$r_i=\sigma(z_i)-\sigma(a_i^\top x^*)$.
По теореме о среднем $r_i=\bar s_i a_i^\top d$, где
$\bar s_i\in[m_\sigma,M_\sigma]$. Поэтому
\[
 \langle\nabla f(x),d\rangle
 =\frac1n\sum_i\bar s_i\sigma'(z_i)(a_i^\top d)^2
 \ge m_\sigma^2d^\top Sd\ge m\|d\|^2.
\]
Та же сумма равна
$n^{-1}\sum_i\sigma'(z_i)r_i^2/\bar s_i
 \ge 2m_\sigma f(x)/M_\sigma$, что доказывает \eqref{reg:quasar}.
Гессиан имеет вид
\[
 \nabla^2f(x)=\frac1n\sum_i
 [\sigma'(z_i)^2+r_i\sigma''(z_i)]a_i a_i^\top.
\]
Для единичного $v$ неравенство Коши-Буняковского и оценка
$(a_i^\top v)^4\le K_A(a_i^\top v)^2$ ограничивают модуль
квадратичной формы величиной
\[
 M_\sigma^2\Lambda+C_\sigma\sqrt{2K_A\Lambda f(x)}
 \le\Lambda\left(M_\sigma^2+\frac{C_\sigma^2}{2t}\right)+K_Atf(x).
\]
Получена \eqref{reg:sighess}. Подстановка
$|r_i|\le\tanh(Z/2)$ даёт \eqref{reg:sigglobal}.
\end{proof}

Приложение включает невыпуклые функции. Возьмём $n=d=1$,
$Q=[-2,2]$, $a_1=1$ и $x^*=-2$. При
$q=\tanh1>1/2$ и $p=\sigma'(2)=(1-q^2)/4$ имеем
\[
 f''(2)=p^2-pq^2=\frac p4(1-5q^2)<0.
\]
По непрерывности кривизна отрицательна и в близких внутренних точках.

Константы требуют точной реализуемости. Для произвольных зашумлённых
ответов эти RSI и связь градиента с ошибкой без остаточного члена
не следуют. При большом $Z$ величина $m_\sigma$ может быть очень
малой, и оценки становятся консервативными. Аффинная граница
кривизны и $L_Q$ обе корректны; ни одна из них не всегда лучше.
Параметр $t$ задаёт границу, а не подбор шага.

Диаметр $\ell_1$-шара равен $D=2R$. Поэтому полный столбцовый ранг
позволяет применять все оценки по норме, зазору и их сочетанию,
включая $\mathcal J_D$; одной \eqref{reg:quasar} достаточно для
остановки с линейной связью ошибки и зазора. Если алгоритму известно,
что ответы реализуемы в допустимой точке, ему уже известно $f^*=0$,
и само значение $f(x)$ равно ошибке. Если при сравнении эта
нижняя граница скрыта, это нужно явно указать. Для произвольной
невыпуклой функции ограниченная задача не объявляется эквивалентной
задаче с фиксированным $\ell_1$-штрафом.
